\documentclass[11pt]{article}
\usepackage[margin=.8in,top=.7in,bottom=.7in]{geometry}
\usepackage{amsmath,amssymb,amsthm,hyperref}
\renewcommand{\normalsize}{\fontsize{10}{12}\selectfont}
\DeclareMathSizes{10}{10}{8}{6}
\hypersetup{colorlinks=true,linkcolor=blue,urlcolor=blue,citecolor=blue}
\setlength{\emergencystretch}{2em}
\newtheorem{theorem}{Theorem}
\newtheorem{lemma}[theorem]{Lemma}
\newcommand{\I}{I_2}
\title{A Hermitian counterexample to\\C*-algebraic Novak positivity}
\author{Research proof and formalization packet}
\date{October 26, 2026}
\begin{document}
\maketitle
\begin{abstract}
Krishna's C*-algebraic Novak conjecture proposes a positive matrix built
from operator cosines of differences of self-adjoint inputs. We give a
counterexample in $M_2(\mathbb C)$ with four inputs and two coordinates.
Every entry of the resulting kernel is a real scalar multiple of the
identity, so the matrix is Hermitian. An explicit vector has quadratic
form $\cos(\pi\sqrt{15})-3\leq-2$. Thus the failure is independent of
product ordering and already occurs before the proposed constant lower
bound is subtracted. A 36-declaration Lean certificate constructs the
literal exponential-defined source matrix and proves its failure of
positive semidefiniteness.
\end{abstract}

\section{The source conjecture}
For an element $T$ of a unital C*-algebra, define
\[
 \cos T=\frac{\exp(iT)+\exp(-iT)}2.
\]
Krishna \cite{K}, Definition~6 and Conjecture~4.2, asks whether
\begin{equation}\label{source}
 \left[\prod_{\ell=1}^{d}
 \frac{1+\cos(X_{j,\ell}-X_{k,\ell})}{2}-\frac1n\right]_{j,k=1}^{n}
 \succeq0
\end{equation}
for all $n,d\geq2$ and all self-adjoint inputs $X_{j,\ell}$.
Here scalar constants denote multiples of the algebra's identity.
The conjecture is numbered~4.3 in the 2021 preprint; the published
Theorem~4.3 proves the commutative case. MathDB~351633 records the
unrestricted conjecture.

\begin{theorem}\label{main}
The assertion \eqref{source} is false for $\mathcal A=M_2(\mathbb C)$,
$n=4$ and $d=2$. It fails for a Hermitian matrix all of whose entries
belong to the center of $\mathcal A$.
\end{theorem}
The original scalar Novak theorem \cite{V} and Krishna's commutative theorem are
unaffected. We use only standard finite-dimensional matrix calculus;
no new general theorem about positive definite radial functions is claimed.

\section{The four inputs}
Let
\[
 Z=\begin{pmatrix}1&0\\0&-1\end{pmatrix},\qquad
 S=\begin{pmatrix}0&1\\1&0\end{pmatrix}.
\]
These matrices satisfy $Z^2=S^2=\I$ and $ZS+SZ=0$. Set
\begin{equation}\label{uv}
 U=2\pi Z,\qquad V=\frac\pi4(-7Z+\sqrt{15}\,S),
 \qquad (P_1,P_2,P_3,P_4)=(0,U,U+V,V).
\end{equation}
Take $X_{j,1}=P_j$ and $X_{j,2}=0$. All eight inputs are self-adjoint.
The anticommutation identity gives
\begin{equation}\label{squares}
 U^2=V^2=4\pi^2\I,\qquad
 (U+V)^2=\pi^2\I,\qquad
 (U-V)^2=15\pi^2\I.
\end{equation}
For example, the $Z$ and $S$ coefficients of $U+V$ are
$\pi/4$ and $\pi\sqrt{15}/4$, whose squared sum is $\pi^2$.

\begin{lemma}\label{cosine}
If $M^2=r^2I$ and $r>0$, then its exponential-defined cosine is
$\cos M=(\cos r)I$.
\end{lemma}
\begin{proof}
Put $E_+=rI+M$ and $E_-=rI-M$. Then
$E_+M=rE_+$ and $E_-M=-rE_-$. Applying the exponential power series
shows
\[
 E_+\cos M=(\cos r)E_+,\qquad
 E_-\cos M=(\cos(-r))E_-=(\cos r)E_-.
\]
Add the identities and use $E_++E_-=2rI$ to cancel the nonzero scalar
$2r$. This argument does not require diagonalization.
\end{proof}
By \eqref{squares}, Lemma~\ref{cosine}, and the evenness of cosine,
\[
 \cos U=\cos V=\I,\qquad
 \cos(U+V)=-\I,\qquad
 \cos(U-V)=c\I,\quad c=\cos(\pi\sqrt{15})\leq1.
\]
Thus, with $t=(1+c)/2$, the uncentered source kernel is
\begin{equation}\label{kernel}
 \left[\prod_{\ell=1}^{2}
 \frac{1+\cos(X_{j,\ell}-X_{k,\ell})}{2}\right]_{j,k=1}^4
 =K\otimes\I,\qquad
 K=\begin{pmatrix}
 1&1&0&1\\1&1&1&t\\0&1&1&1\\1&t&1&1
 \end{pmatrix}.
\end{equation}
Indeed, the second-coordinate factor is always $\I$.
All factors appearing here are scalar matrices, so either ordered or
symmetrized multiplication gives exactly \eqref{kernel}.

\begin{proof}[Proof of Theorem~\ref{main}]
Let $y=(1,-1,1,-1)^T$. Since $\sum_j y_j=0$, the constant matrix
$J_4$ satisfies $y^TJ_4y=0$. Direct multiplication yields
\[
 y^T(K-J_4/4)y=y^TKy=-4+2t=c-3\leq-2<0.
\]
Consequently the literal matrix in \eqref{source}, which is
$(K-J_4/4)\otimes\I$, has a negative quadratic form on
$y\otimes(1,0)^T$. It is Hermitian but not positive semidefinite.
The same computation also rules out positivity of the uncentered kernel.
\end{proof}

\section{Formal scope and provenance}
The three-module Lean certificate has 36 declarations. It constructs the
operator cosine from \texttt{NormedSpace.exp}, proves all input and cosine
identities, and defines the literal source matrix. Its endpoint,\\
\texttt{Novak.source\_counterexample\_351633}, proves self-adjoint inputs,
a Hermitian source matrix, and failure of positive semidefiniteness,
without residual hypotheses. Every declaration passes an isolated strict
replay with Lean~4.34.1 and pinned mathlib revision\\
\texttt{d13f23b723b8a846827a245b89c10fc7d3f11612}; only
\texttt{propext}, \texttt{Classical.choice} and \texttt{Quot.sound} occur
as logical axioms.

The published source statement and cosine definition were checked directly
on page~799. Source and literature screening, independent review and
artifact receipts are recorded separately in the accompanying research
packet. No absolute priority or separate discovery credit for helpers is
claimed.

\begin{thebibliography}{9}
\bibitem{K} K.~M.~Krishna,
C*-algebraic Schur product theorem, P\'olya--Szeg\H{o}--Rudin question
and Novak's conjecture,
\emph{J. Korean Math. Soc.} \textbf{59} (2022), no.~4, 789--804.
\href{https://doi.org/10.4134/JKMS.j210627}{doi:10.4134/JKMS.j210627}.
Preprint: \href{https://arxiv.org/abs/2108.06662}{arXiv:2108.06662}.
\bibitem{V} J.~Vyb\'iral,
A variant of Schur's product theorem and its applications,
\emph{Adv. Math.} \textbf{368} (2020), 107140.
\href{https://doi.org/10.1016/j.aim.2020.107140}{doi:10.1016/j.aim.2020.107140}.
\end{thebibliography}
\end{document}
